\begin{table}[H]
\centering
\scriptsize
\renewcommand{\arraystretch}{1.15}
\setlength{\tabcolsep}{6pt}
\begin{tabular}{|l|r|r|r|r|}
\hline
\multicolumn{5}{|c|}{\textbf{Desarrollo de Servicio}} \\
\multicolumn{5}{|c|}{\textbf{Costo de capital}} \\ \hline

\textbf{Concepto} & \textbf{Monto de la Inversión} & \textbf{\%} & \textbf{Costo} & \textbf{Costo Promedio} \\ \hline

Fondos propios & 1,154.29 & 6.0263\% & 6.000\% & 0.3616\% \\ \hline
Financiamiento & 18,000.00 & 93.9737\% & 12.000\% & 11.2768\% \\ \hline
Costo de capital &  &  &  & 11.6384\% \\ \hline
Inflación promedio &  &  &  & 2.2900\% \\ \hline
Riesgo de la inversión &  &  &  & 0.2665\% \\ \hline
\textbf{Tasa de Corte} &  &  &  & \textbf{14.194942\%} \\ \hline

\end{tabular}

\vspace{0.2cm}
\textbf{Figura 10.28.} Costo de Capital \textbf{[Elaboración Propia]}
\end{table}

\textbf{Valor Actual Neto VAN}

El valor actual neto responde a la siguiente fórmula en Excel =VAN(tasa de corte, flujo año1, flujo año2, flujo año3, flujo año4, flujo año5, flujo año6).

\textbf{Periodo de Recuperación}

El análisis del periodo de recuperación utilizando flujos descontados permite determinar en qué momento el proyecto logra recuperar la inversión inicial considerando el valor del dinero en el tiempo. En este caso, la inversión inicial asciende a S/ 21,154.29. Al descontar los flujos de caja netos proyectados para cada año, se observa que el acumulado aún es negativo al finalizar el segundo año. No obstante, en el tercer año, el flujo descontado (S/ 99,829.15) permite cubrir el saldo pendiente, alcanzando un valor acumulado positivo. Mediante interpolación, se estima que el punto de equilibrio ocurre a los 2.79 años. Esto significa que el proyecto recupera su inversión inicial a través de los flujos netos descontados aproximadamente entre el segundo y tercer año de operación, lo cual refleja una recuperación relativamente rápida considerando la naturaleza tecnológica del servicio.